\chapter{Complexity — DAG Circuit}
%%% generated by the blueprint pipeline from TCSlib.Complexity.CircuitComplexity.DAGCircuit
\section{Overview}

Declarations of the Lean module \texttt{TCSlib.Complexity.CircuitComplexity.DAGCircuit}.

\section{Declarations}

\begin{definition}[Boolean gate kinds]
\lean{BoolCircuit.GateKind}
\label{BoolCircuit.GateKind}
\leanok
A Boolean gate is an AND gate, an OR gate, or a NOT gate.
\end{definition}

\begin{definition}[Gate labels and inputs]
\lean{BoolCircuit.DAGGate}
\label{BoolCircuit.DAGGate}
\leanok
\uses{BoolCircuit.GateKind}
A directed acyclic graph gate consists of a gate kind and an ordered list of
natural-number indices identifying the values it reads.
\end{definition}

\begin{definition}[Gate evaluation]
\lean{BoolCircuit.DAGGate.eval}
\label{BoolCircuit.DAGGate.eval}
\leanok
\uses{BoolCircuit.DAGGate}
Given a gate and a list of Boolean values, evaluate an AND gate by taking the
conjunction of the values at its listed indices, an OR gate by taking their disjunction,
and a NOT gate by negating their conjunction. Any listed index outside the value list
contributes false.
\end{definition}

\begin{definition}[Gate depth]
\lean{BoolCircuit.DAGGate.depth}
\label{BoolCircuit.DAGGate.depth}
\leanok
\uses{BoolCircuit.DAGGate}
Given a gate and a list of natural-number depths, its depth is one plus the maximum of
zero and the depths at its listed indices; an index outside the depth list contributes
zero.
\end{definition}

\begin{definition}[Sequential gate execution]
\lean{BoolCircuit.runWith}
\label{BoolCircuit.runWith}
\leanok
\uses{BoolCircuit.DAGGate}
Given a function that computes a value from a gate and the list of values already
available, and an initial list of values, process the gates in order, appending each
computed value to the list, and return the resulting list.
\end{definition}

\begin{theorem}[Running an empty gate list]
\lean{BoolCircuit.runWith_nil}
\label{BoolCircuit.runWith_nil}
\leanok
\uses{BoolCircuit.runWith}
\difficulty{1}
Let $\beta$ be a type, let $f$ assign an output in $\beta$ to each gate and list of
previous outputs, and let $init$ be an initial output list. Running an empty gate list
from $init$ returns $init$.
\end{theorem}

\begin{theorem}[Execution beginning with one gate]
\lean{BoolCircuit.runWith_cons}
\label{BoolCircuit.runWith_cons}
\leanok
\uses{BoolCircuit.DAGGate, BoolCircuit.runWith}
\difficulty{1}
Let $\beta$ be a type, let $f$ assign a value in $\beta$ to each gate and current list
of values, let $g$ be a gate, let $g_s$ be a list of gates, and let $I$ be an initial
list of values. Executing $g::g_s$ from $I$ gives the same list as executing $g_s$ from
the list obtained by appending $f(g,I)$ to $I$.
\end{theorem}

\begin{theorem}[Execution over concatenated gate lists]
\lean{BoolCircuit.runWith_append}
\label{BoolCircuit.runWith_append}
\leanok
\uses{BoolCircuit.DAGGate, BoolCircuit.runWith}
\difficulty{3}
Let $\beta$ be a type, let $f$ assign a value in $\beta$ to each gate and current list
of values, and let $g_s,h_s$ be lists of gates and $I$ an initial list of values.
Executing the concatenated gate list $g_s++h_s$ from $I$ gives the same list as first
executing $g_s$ from $I$ and then executing $h_s$ from the resulting list.
\end{theorem}

\begin{theorem}[Running a single gate]
\lean{BoolCircuit.runWith_singleton}
\label{BoolCircuit.runWith_singleton}
\leanok
\uses{BoolCircuit.DAGGate, BoolCircuit.runWith}
\difficulty{1}
Let $\beta$ be a type, let $f$ assign an output in $\beta$ to each gate and list of
previous outputs, let $g$ be a gate, and let $init$ be an initial output list. Running
the one-gate list $[g]$ from $init$ appends the output $f(g,init)$ to $init$.
\end{theorem}

\begin{theorem}[Length after evaluating gates]
\lean{BoolCircuit.length_runWith}
\label{BoolCircuit.length_runWith}
\leanok
\uses{BoolCircuit.DAGGate, BoolCircuit.runWith, BoolCircuit.runWith_cons}
\difficulty{4}
For any function $f$ assigning a value to a gate and a current list of values, any
finite list $gs$ of gates and initial list $init$ produce a list of length $|init|+|gs|$
when the gates are run in sequence.
\end{theorem}

\begin{theorem}[Appended gate outputs]
\lean{BoolCircuit.runWith_eq_append}
\label{BoolCircuit.runWith_eq_append}
\leanok
\uses{BoolCircuit.DAGGate, BoolCircuit.runWith, BoolCircuit.runWith_cons}
\difficulty{4}
For any type of values, gate-evaluation function, finite list of gates $g_1,\ldots,g_m$,
and initial value list $v$, successively append to the current list the value returned
by evaluating each gate on the current list. The resulting list is $v$ followed by some
list of values.
\end{theorem}

\begin{theorem}[Initial values are preserved]
\lean{BoolCircuit.runWith_getD_of_lt}
\label{BoolCircuit.runWith_getD_of_lt}
\leanok
\uses{BoolCircuit.DAGGate, BoolCircuit.runWith, BoolCircuit.runWith_eq_append}
\difficulty{3}
Let $f$ assign a value to each gate and current value list, let $G$ be a finite gate
list, and let $v$ be an initial value list. For any index $j<|v|$ and default value $d$,
the value at index $j$ after evaluating $G$ from $v$ is the value at index $j$ of $v$,
using $d$ if that index is absent.
\end{theorem}

\begin{theorem}[Value produced by the last gate]
\lean{BoolCircuit.runWith_getD_last}
\label{BoolCircuit.runWith_getD_last}
\leanok
\uses{BoolCircuit.DAGGate, BoolCircuit.length_runWith, BoolCircuit.runWith, BoolCircuit.runWith_append, BoolCircuit.runWith_singleton}
\difficulty{3}
Let $f$ assign a value to each gate and current value list, let $g$ be a gate, and let
$v$ be an initial value list. For any finite gate list $G$ and default value $d$, the
value at index $|v|+|G|$ after evaluating $G$ followed by $g$ is $f(g,w)$, where $w$ is
the list obtained by evaluating $G$ starting from $v$.
\end{theorem}

\begin{theorem}[Value produced by a gate]
\lean{BoolCircuit.runWith_getD_gate}
\label{BoolCircuit.runWith_getD_gate}
\leanok
\uses{BoolCircuit.DAGGate, BoolCircuit.runWith, BoolCircuit.runWith_append, BoolCircuit.runWith_getD_last, BoolCircuit.runWith_getD_of_lt}
\difficulty{4}
Let $f$ assign a value to each gate and current value list, and let $g_0,\ldots,g_{m-1}$
be a finite gate list with initial value list $v$. For any index $i<m$ and default value
$d$, the value at index $|v|+i$ after successively evaluating the gates is $f(g_i, w)$,
where $w$ is the list obtained by evaluating just the first $i$ gates starting from $v$.
\end{theorem}

\begin{theorem}[Prefix evaluation preserves earlier values]
\lean{BoolCircuit.runWith_getD_take}
\label{BoolCircuit.runWith_getD_take}
\leanok
\uses{BoolCircuit.DAGGate, BoolCircuit.length_runWith, BoolCircuit.runWith, BoolCircuit.runWith_append, BoolCircuit.runWith_getD_of_lt}
\difficulty{3}
Let $f$ assign a value to each gate and current value list, let $G$ be a finite gate
list, and let $v$ be an initial value list. If $i\le |G|$ and $j<|v|+i$, then the value
at index $j$ after evaluating the first $i$ gates from $v$ agrees with the value at that
index after evaluating all of $G$, using any default value $d$ when an index is absent.
\end{theorem}

\begin{theorem}[Equality of list-wide Boolean tests]
\lean{BoolCircuit.all_congr_mem}
\label{BoolCircuit.all_congr_mem}
\leanok
\difficulty{4}
Let $l$ be a list in a type $\alpha$, and let
$p,q:\alpha\to\{\mathrm{false},\mathrm{true}\}$. If $p(a)=q(a)$ for every entry $a$ of
$l$, then the Boolean test that every entry of $l$ satisfies $p$ equals the
corresponding test for $q$.
\end{theorem}

\begin{theorem}[Congruence of list disjunction]
\lean{BoolCircuit.any_congr_mem}
\label{BoolCircuit.any_congr_mem}
\leanok
\difficulty{4}
Let $l$ be a list and let $p$ and $q$ be Boolean-valued predicates on its elements. If
$p(a)=q(a)$ for every element $a$ of $l$, then the listwise disjunction of $p$ over $l$
equals the listwise disjunction of $q$ over $l$.
\end{theorem}

\begin{theorem}[Evaluation depends on referenced values]
\lean{BoolCircuit.DAGGate.eval_congr}
\label{BoolCircuit.DAGGate.eval_congr}
\leanok
\uses{BoolCircuit.DAGGate, BoolCircuit.DAGGate.eval, BoolCircuit.all_congr_mem, BoolCircuit.any_congr_mem}
\difficulty{4}
Let $g$ be a Boolean circuit gate and let $vals,vals'$ be lists of Boolean values. If
their defaulted lookups agree at every argument of $g$, using false when an index is out
of range, then $g$ evaluates to the same Boolean value on both lists.
\end{theorem}

\begin{theorem}[Depth depends on referenced values]
\lean{BoolCircuit.DAGGate.depth_congr}
\label{BoolCircuit.DAGGate.depth_congr}
\leanok
\uses{BoolCircuit.DAGGate, BoolCircuit.DAGGate.depth}
\difficulty{3}
Let $g$ be a Boolean circuit gate and let $ds,ds'$ be lists of natural numbers. If the
values looked up in these lists at every argument of $g$ agree, using $0$ when an index
is out of range, then the depth of $g$ computed from $ds$ equals its depth computed from
$ds'$.
\end{theorem}

\begin{definition}[Well-formed gate]
\lean{BoolCircuit.DAGGate.WellFormed}
\label{BoolCircuit.DAGGate.WellFormed}
\leanok
\uses{BoolCircuit.DAGGate}
A gate is well formed when its list of input indices has no repetitions and, if it is a
NOT gate, that list has length exactly one.
\end{definition}

\begin{definition}[Gate with at most two inputs]
\lean{BoolCircuit.DAGGate.FaninTwo}
\label{BoolCircuit.DAGGate.FaninTwo}
\leanok
\uses{BoolCircuit.DAGGate, BoolCircuit.DAGGate.WellFormed}
A gate has fan-in at most two when it is well formed and its list of input indices has
length at most two.
\end{definition}

\begin{definition}[Boolean constant gate]
\lean{BoolCircuit.constGate}
\label{BoolCircuit.constGate}
\leanok
\uses{BoolCircuit.DAGGate}
For a Boolean value $b$, the constant gate is an AND gate with no inputs when $b$ is
true, and an OR gate with no inputs when $b$ is false.
\end{definition}

\begin{theorem}[True Constant Gate]
\lean{BoolCircuit.constGate_true}
\label{BoolCircuit.constGate_true}
\leanok
\uses{BoolCircuit.constGate}
\difficulty{1}
A constant gate labelled by true is the AND gate with no input vertices.
\end{theorem}

\begin{theorem}[False Constant Gate]
\lean{BoolCircuit.constGate_false}
\label{BoolCircuit.constGate_false}
\leanok
\uses{BoolCircuit.constGate}
\difficulty{1}
A constant gate labelled by false is the OR gate with no input vertices.
\end{theorem}

\begin{theorem}[Evaluation of a constant gate]
\lean{BoolCircuit.constGate_eval}
\label{BoolCircuit.constGate_eval}
\leanok
\uses{BoolCircuit.DAGGate.eval, BoolCircuit.constGate}
\difficulty{4}
For every Boolean value $b$ and every finite list of Boolean values, the gate
representing the constant $b$ evaluates to $b$, regardless of the list.
\end{theorem}

\begin{theorem}[Depth of a constant gate]
\lean{BoolCircuit.constGate_depth}
\label{BoolCircuit.constGate_depth}
\leanok
\uses{BoolCircuit.DAGGate.depth, BoolCircuit.constGate}
\difficulty{4}
For every Boolean value $b$ and every list of natural-number depths assigned to
vertices, the gate representing the constant $b$ has depth $1$.
\end{theorem}

\begin{theorem}[Arguments of a constant gate]
\lean{BoolCircuit.constGate_args}
\label{BoolCircuit.constGate_args}
\leanok
\uses{BoolCircuit.constGate}
\difficulty{1}
For every Boolean value $b$, the gate representing the constant $b$ has no input
arguments.
\end{theorem}

\begin{theorem}[Constant Gate Is Not a NOT Gate]
\lean{BoolCircuit.constGate_kind_ne_not}
\label{BoolCircuit.constGate_kind_ne_not}
\leanok
\uses{BoolCircuit.constGate}
\difficulty{3}
For either Boolean value, the corresponding constant gate is not labelled by NOT.
\end{theorem}

\begin{theorem}[Constant Gate Fan-In Bound]
\lean{BoolCircuit.constGate_faninTwo}
\label{BoolCircuit.constGate_faninTwo}
\leanok
\uses{BoolCircuit.DAGGate.FaninTwo, BoolCircuit.constGate, BoolCircuit.constGate_kind_ne_not}
\difficulty{3}
For either Boolean value, the corresponding constant gate is well formed and has at most
two input vertices.
\end{theorem}

\begin{definition}[Gate argument remapping]
\lean{BoolCircuit.DAGGate.remap}
\label{BoolCircuit.DAGGate.remap}
\leanok
\uses{BoolCircuit.DAGGate}
Given a map from natural numbers to natural numbers and a gate, form a gate of the same
kind whose argument list is obtained by applying the map to each entry of the original
list.
\end{definition}

\begin{theorem}[Identity remapping of a gate]
\lean{BoolCircuit.DAGGate.remap_id}
\label{BoolCircuit.DAGGate.remap_id}
\leanok
\uses{BoolCircuit.DAGGate, BoolCircuit.DAGGate.remap}
\difficulty{3}
For every Boolean circuit gate $g$, remapping its argument indices by the identity
function leaves $g$ unchanged.
\end{theorem}

\begin{theorem}[Evaluation under input remapping]
\lean{BoolCircuit.DAGGate.eval_remap}
\label{BoolCircuit.DAGGate.eval_remap}
\leanok
\uses{BoolCircuit.DAGGate, BoolCircuit.DAGGate.eval, BoolCircuit.DAGGate.remap, BoolCircuit.all_congr_mem, BoolCircuit.any_congr_mem}
\difficulty{4}
Let $g$ be a gate, let $\sigma$ map natural-number indices to natural-number indices,
and let $v,v'$ be Boolean value lists. If for every input index $a$ of $g$, the value of
$v$ at index $\sigma(a)$ (with out-of-range indices treated as false) equals the value
of $v'$ at index $a$, then the gate obtained by replacing each input index $a$ with
$\sigma(a)$ has the same evaluation on $v$ as $g$ has on $v'$.
\end{theorem}

\begin{theorem}[Relation preserved under gate remapping]
\lean{BoolCircuit.runWith_remap_rel}
\label{BoolCircuit.runWith_remap_rel}
\leanok
\uses{BoolCircuit.DAGGate, BoolCircuit.DAGGate.remap, BoolCircuit.runWith, BoolCircuit.runWith_append, BoolCircuit.runWith_singleton}
\difficulty{7}
Let $\beta$ be a type, let $f$ assign an output in $\beta$ to each gate and list of
previous outputs, let $R$ be a relation on $\beta$, let $d\in\beta$, and let
$\sigma:\mathbb{N}\to\mathbb{N}$. Suppose that for every gate $g$ and output lists
$v,v'$, if $R(v'_{\sigma(a)},v_a)$ holds for every argument index $a$ of $g$ (where an
out-of-range list index returns $d$), then $R(f(g\text{ with every argument index
}a\text{ replaced by }\sigma(a),v),f(g,v'))$ holds. Let $L,L'\in\mathbb{N}$ and let
$init,init'$ be lists of lengths $L,L'$, respectively. Suppose $\sigma(v)<L'$ whenever
$v<L$, $\sigma(L+i)=L'+i$ for every $i$, and $R(init'_{\sigma(v)},init_v)$ for every
$v<L$, again using $d$ for out-of-range indices. Let $gs$ be a gate list such that, for
each position $i$ in $gs$, every argument index of its gate is less than $L+i$. Then for
every $v<L+|gs|$, the output at index $\sigma(v)$ from running the remapped gates on
$init'$ is related by $R$ to the output at index $v$ from running $gs$ on $init$, with
out-of-range indices interpreted as $d$.
\end{theorem}

\begin{definition}[Acyclicity condition for gates]
\lean{BoolCircuit.GatesAcyclic}
\label{BoolCircuit.GatesAcyclic}
\leanok
\uses{BoolCircuit.DAGGate}
For a natural number $n$ and a finite ordered list of gates, every argument of the gate
at position $i$ is less than $n+i$, for each valid position $i$ in the list.
\end{definition}

\begin{theorem}[Acyclicity of the empty gate list]
\lean{BoolCircuit.GatesAcyclic.nil}
\label{BoolCircuit.GatesAcyclic.nil}
\leanok
\uses{BoolCircuit.GatesAcyclic}
\difficulty{1}
For every natural number $n$, the empty list of gates satisfies the acyclicity condition
starting at $n$.
\end{theorem}

\begin{theorem}[Acyclicity under appending a gate]
\lean{BoolCircuit.GatesAcyclic.snoc}
\label{BoolCircuit.GatesAcyclic.snoc}
\leanok
\uses{BoolCircuit.DAGGate, BoolCircuit.GatesAcyclic}
\difficulty{4}
Let $n$ be the number of input vertices and let $g_0,\ldots,g_{r-1}$ be an ordered list
of gates such that every source of gate $g_i$ has index less than $n+i$. If every source
of a further gate $g$ has index less than $n+r$, then the list obtained by appending $g$
has the same property.
\end{theorem}

\begin{theorem}[Acyclicity under concatenation]
\lean{BoolCircuit.GatesAcyclic.append}
\label{BoolCircuit.GatesAcyclic.append}
\leanok
\uses{BoolCircuit.DAGGate, BoolCircuit.GatesAcyclic}
\difficulty{4}
Let $n$ be a natural number and let $gs$ and $hs$ be lists of gates. If every gate at
position $i$ in $gs$ refers only to indices less than $n+i$, and every gate at position
$i$ in $hs$ refers only to indices less than $n+|gs|+i$, then every gate at position $i$
in the concatenated list $gs$ followed by $hs$ refers only to indices less than $n+i$.
\end{theorem}

\begin{theorem}[Acyclicity for the empty gate list]
\lean{BoolCircuit.gatesAcyclic_nil}
\label{BoolCircuit.gatesAcyclic_nil}
\leanok
\uses{BoolCircuit.GatesAcyclic, BoolCircuit.GatesAcyclic.nil}
\difficulty{1}
For every $n\in\mathbb{N}$, the empty list of gates is acyclic relative to $n$ if and
only if true.
\end{theorem}

\begin{theorem}[Acyclicity of a list with a first gate]
\lean{BoolCircuit.gatesAcyclic_cons}
\label{BoolCircuit.gatesAcyclic_cons}
\leanok
\uses{BoolCircuit.DAGGate, BoolCircuit.GatesAcyclic}
\difficulty{4}
Let $n$ be a natural number, let $g$ be a gate, and let $gs$ be a finite list of gates.
The list consisting of $g$ followed by $gs$ is acyclic starting from index bound $n$ if
and only if every input index read by $g$ is less than $n$ and $gs$ is acyclic starting
from $n+1$.
\end{theorem}

\begin{theorem}[Acyclicity under concatenation]
\lean{BoolCircuit.gatesAcyclic_append}
\label{BoolCircuit.gatesAcyclic_append}
\leanok
\uses{BoolCircuit.DAGGate, BoolCircuit.GatesAcyclic, BoolCircuit.gatesAcyclic_cons}
\difficulty{4}
Let $n$ be a natural number and let $gs$ and $hs$ be finite lists of gates. Their
concatenation is acyclic starting from index bound $n$ if and only if $gs$ is acyclic
starting from $n$ and $hs$ is acyclic starting from $n+|gs|$, where $|gs|$ is the number
of gates in $gs$.
\end{theorem}

\begin{definition}[Directed acyclic circuit]
\lean{BoolCircuit.DAGCircuit}
\label{BoolCircuit.DAGCircuit}
\leanok
\uses{BoolCircuit.DAGGate}
A circuit with $n$ input vertices consists of a finite ordered list of gates and an
output vertex. Each gate at position $i$ reads only vertices with index less than $n+i$,
and the output vertex has index less than $n$ plus the number of gates.
\end{definition}

\begin{definition}[Circuit value sequence]
\lean{BoolCircuit.DAGCircuit.values}
\label{BoolCircuit.DAGCircuit.values}
\leanok
\uses{BoolCircuit.DAGGate.eval, BoolCircuit.runWith}
Given a circuit with $n$ input vertices and an assignment of Boolean values to those
inputs, the value sequence is obtained by running the circuit's gates, in order, on the
list of input values.
\end{definition}

\begin{definition}[Circuit evaluation]
\lean{BoolCircuit.DAGCircuit.eval}
\label{BoolCircuit.DAGCircuit.eval}
\leanok
\uses{BoolCircuit.DAGCircuit.values}
Given a circuit with $n$ input vertices and an assignment of Boolean values to those
inputs, its output is the value at the circuit's designated output index in the value
sequence; if that index is out of range, the output is false.
\end{definition}

\begin{definition}[Circuit vertex depths]
\lean{BoolCircuit.DAGCircuit.depths}
\label{BoolCircuit.DAGCircuit.depths}
\leanok
\uses{BoolCircuit.DAGGate.depth, BoolCircuit.runWith}
For a circuit with $n$ input vertices, the depth list is obtained by running the
circuit's gates, in order, starting from a list of $n$ zeros.
\end{definition}

\begin{definition}[Depth at a vertex]
\lean{BoolCircuit.DAGCircuit.depthAt}
\label{BoolCircuit.DAGCircuit.depthAt}
\leanok
\uses{BoolCircuit.DAGCircuit.depths}
For a circuit with $n$ input vertices and a natural-number vertex index $v$, the depth
at $v$ is the entry of the circuit's depth list at index $v$, or zero if that index is
out of range.
\end{definition}

\begin{definition}[Circuit depth]
\lean{BoolCircuit.DAGCircuit.depth}
\label{BoolCircuit.DAGCircuit.depth}
\leanok
\uses{BoolCircuit.DAGCircuit.depthAt}
For a circuit with $n$ input vertices, its depth is the depth at its designated output
vertex.
\end{definition}

\begin{definition}[Circuit size]
\lean{BoolCircuit.DAGCircuit.size}
\label{BoolCircuit.DAGCircuit.size}
\leanok
For a directed acyclic circuit with $n$ input vertices and a list of gates, define its
size to be $n$ plus the number of gates.
\end{definition}

\begin{definition}[Well-formed circuit]
\lean{BoolCircuit.DAGCircuit.IsWellFormed}
\label{BoolCircuit.DAGCircuit.IsWellFormed}
\leanok
A directed acyclic circuit is well formed if every gate reads distinct vertices, and
every NOT gate reads exactly one vertex.
\end{definition}

\begin{definition}[Circuit with fan-in at most two]
\lean{BoolCircuit.DAGCircuit.IsFaninTwo}
\label{BoolCircuit.DAGCircuit.IsFaninTwo}
\leanok
\uses{BoolCircuit.DAGCircuit.IsWellFormed}
A directed acyclic circuit has fan-in at most two if it is well formed and every gate
reads at most two vertices.
\end{definition}

\begin{theorem}[Length of the circuit value list]
\lean{BoolCircuit.DAGCircuit.length_values}
\label{BoolCircuit.DAGCircuit.length_values}
\leanok
\uses{BoolCircuit.DAGCircuit.values}
\difficulty{1}
For a directed acyclic circuit with $n$ input vertices and a list of gates in
topological order, and for any Boolean assignment to its inputs, the list of values
obtained by starting with the input values and evaluating the gates in order has length
$n$ plus the number of gates.
\end{theorem}

\begin{theorem}[Length of the depth list]
\lean{BoolCircuit.DAGCircuit.length_depths}
\label{BoolCircuit.DAGCircuit.length_depths}
\leanok
\uses{BoolCircuit.DAGCircuit.depths}
\difficulty{1}
For a directed acyclic circuit with $n$ input vertices and a list of gates in
topological order, the list of vertex depths has length $n$ plus the number of gates.
\end{theorem}

\begin{theorem}[Computed values at input vertices]
\lean{BoolCircuit.DAGCircuit.values_getD_input}
\label{BoolCircuit.DAGCircuit.values_getD_input}
\leanok
\uses{BoolCircuit.DAGCircuit.values, BoolCircuit.runWith_getD_of_lt}
\difficulty{3}
Let $C$ be a Boolean directed acyclic circuit with $n$ inputs, and let $x$ be an
assignment to those inputs. At every input vertex $i$, the computed value is $x(i)$.
\end{theorem}

\begin{theorem}[Depth at an input]
\lean{BoolCircuit.DAGCircuit.depthAt_input}
\label{BoolCircuit.DAGCircuit.depthAt_input}
\leanok
\uses{BoolCircuit.DAGCircuit.depthAt, BoolCircuit.DAGCircuit.depths, BoolCircuit.runWith_getD_of_lt}
\difficulty{3}
Let $C$ be a directed acyclic Boolean circuit on $n$ input vertices. Every input vertex
$i<n$ has depth zero.
\end{theorem}

\begin{theorem}[Computed value at a gate vertex]
\lean{BoolCircuit.DAGCircuit.values_getD_gate}
\label{BoolCircuit.DAGCircuit.values_getD_gate}
\leanok
\uses{BoolCircuit.DAGCircuit.values, BoolCircuit.DAGGate.eval, BoolCircuit.DAGGate.eval_congr, BoolCircuit.runWith_getD_gate, BoolCircuit.runWith_getD_take}
\difficulty{4}
Let $C$ be a Boolean directed acyclic circuit with $n$ inputs, and let $x$ be an
assignment to its inputs. For every gate index $i<|C.\text{gates}|$, the computed value
at vertex $n+i$ equals the evaluation of gate $i$ on the list of values computed by the
circuit.
\end{theorem}

\begin{theorem}[Depth at a gate]
\lean{BoolCircuit.DAGCircuit.depthAt_gate}
\label{BoolCircuit.DAGCircuit.depthAt_gate}
\leanok
\uses{BoolCircuit.DAGCircuit.depthAt, BoolCircuit.DAGCircuit.depths, BoolCircuit.DAGGate.depth, BoolCircuit.DAGGate.depth_congr, BoolCircuit.runWith_getD_gate, BoolCircuit.runWith_getD_take}
\difficulty{4}
Let $C$ be a directed acyclic Boolean circuit on $n$ input vertices, and let $i$ be a
gate index with $i<|C|$. The depth assigned to vertex $n+i$ equals one plus the maximum
of zero and the depths of the vertices read by gate $i$, where any depth lookup beyond
the available depth list is taken to be zero.
\end{theorem}

\begin{definition}[Constant circuit]
\lean{BoolCircuit.constCircuit}
\label{BoolCircuit.constCircuit}
\leanok
\uses{BoolCircuit.DAGCircuit, BoolCircuit.constGate}
For a Boolean value $b$ and a number $n$ of inputs, the constant circuit consists of one
gate representing $b$ and has output vertex $n$.
\end{definition}

\begin{theorem}[Evaluation of a constant circuit]
\lean{BoolCircuit.constCircuit_eval}
\label{BoolCircuit.constCircuit_eval}
\leanok
\uses{BoolCircuit.DAGCircuit.eval, BoolCircuit.DAGCircuit.values, BoolCircuit.constCircuit, BoolCircuit.runWith_cons}
\difficulty{1}
For every natural number $n$, Boolean value $b$, and input $x\in\{0,1\}^n$, the circuit
with one gate representing the constant $b$ evaluates to $b$ on $x$.
\end{theorem}

\begin{theorem}[Constant circuits have fan-in two]
\lean{BoolCircuit.constCircuit_isFaninTwo}
\label{BoolCircuit.constCircuit_isFaninTwo}
\leanok
\uses{BoolCircuit.DAGCircuit.IsFaninTwo, BoolCircuit.constCircuit, BoolCircuit.constGate}
\difficulty{4}
For every natural number $n$ and Boolean value $b$, the circuit with one gate
representing the constant $b$ is well formed and every gate has at most two inputs.
\end{theorem}

\begin{theorem}[Size of a constant circuit]
\lean{BoolCircuit.constCircuit_size}
\label{BoolCircuit.constCircuit_size}
\leanok
\uses{BoolCircuit.DAGCircuit.size, BoolCircuit.constCircuit}
\difficulty{1}
For every natural number $n$ and Boolean value $b$, the circuit with one gate
representing the constant $b$ has size $n+1$.
\end{theorem}

\begin{definition}[Family of circuits]
\lean{BoolCircuit.DAGCircuitFamily}
\label{BoolCircuit.DAGCircuitFamily}
\leanok
\uses{BoolCircuit.DAGCircuit}
A family of directed acyclic circuits assigns to each natural number $n$ a circuit with
$n$ input vertices.
\end{definition}

\begin{definition}[Acceptance by a circuit family]
\lean{BoolCircuit.DAGCircuitFamily.Accepts}
\label{BoolCircuit.DAGCircuitFamily.Accepts}
\leanok
\uses{BoolCircuit.DAGCircuit.eval}
A circuit family accepts a Boolean word $w$ if the circuit assigned to the length of
$w$, evaluated on the word's entries, outputs true.
\end{definition}

\begin{definition}[Language of a circuit family]
\lean{BoolCircuit.DAGCircuitFamily.language}
\label{BoolCircuit.DAGCircuitFamily.language}
\leanok
\uses{BoolCircuit.DAGCircuitFamily.Accepts}
The language of a circuit family is the set of Boolean words that it accepts.
\end{definition}

\begin{theorem}[Membership in a circuit family's language]
\lean{BoolCircuit.DAGCircuitFamily.mem_language_iff}
\label{BoolCircuit.DAGCircuitFamily.mem_language_iff}
\leanok
\uses{BoolCircuit.DAGCircuit.eval, BoolCircuit.DAGCircuitFamily.language}
\difficulty{1}
Let $C$ be a family of Boolean circuits and let $w$ be a finite Boolean string. Then $w$
belongs to the language accepted by $C$ if and only if the circuit in the family for
inputs of length $|w|$ outputs true on $w$.
\end{theorem}

\begin{definition}[Well-formed circuit family]
\lean{BoolCircuit.DAGCircuitFamily.IsWellFormed}
\label{BoolCircuit.DAGCircuitFamily.IsWellFormed}
\leanok
\uses{BoolCircuit.DAGCircuit.IsWellFormed}
A circuit family is well formed if, for every natural number $n$, its circuit for inputs
of length $n$ is well formed.
\end{definition}

\begin{definition}[Fan-in-two circuit family]
\lean{BoolCircuit.DAGCircuitFamily.HasFaninTwo}
\label{BoolCircuit.DAGCircuitFamily.HasFaninTwo}
\leanok
\uses{BoolCircuit.DAGCircuit.IsFaninTwo}
A family of directed acyclic circuits has fan-in two if, for every input size $n$, every
gate in its circuit on $n$ inputs has at most two incoming arguments.
\end{definition}

\begin{definition}[Polynomial-size circuit family]
\lean{BoolCircuit.DAGCircuitFamily.IsPolySize}
\label{BoolCircuit.DAGCircuitFamily.IsPolySize}
\leanok
\uses{BoolCircuit.DAGCircuit.size}
A family of directed acyclic circuits has polynomial size if there are natural numbers
$a$ and $k$ such that, for every input size $n$, the circuit has at most $a(n+1)^k$
gates.
\end{definition}

\begin{definition}[Polylogarithmic circuit depth]
\lean{BoolCircuit.DAGCircuitFamily.HasPolylogDepth}
\label{BoolCircuit.DAGCircuitFamily.HasPolylogDepth}
\leanok
\uses{BoolCircuit.DAGCircuit.depth}
For a natural number $d$, a family of directed acyclic circuits has depth bounded by a
degree-$d$ polylogarithm if there is a natural number $b$ such that, for every input
size $n$, its circuit has depth at most $b(loor{\text{log}_2 n}+1)^d$.
\end{definition}

\begin{theorem}[Fanin two implies well formedness]
\lean{BoolCircuit.DAGCircuitFamily.HasFaninTwo.isWellFormed}
\label{BoolCircuit.DAGCircuitFamily.HasFaninTwo.isWellFormed}
\leanok
\uses{BoolCircuit.DAGCircuitFamily, BoolCircuit.DAGCircuitFamily.HasFaninTwo, BoolCircuit.DAGCircuitFamily.IsWellFormed}
\difficulty{1}
Let $C$ be a family of Boolean directed acyclic circuits, with one circuit for each
input length. If every circuit in the family is well formed and every gate in each
circuit has at most two arguments, then every circuit in the family is well formed.
\end{theorem}
